\chapter{\rm\bfseries Experimental Design and Procedure}
\label{ch:experiments}

The empirical campaign is organised into two sequential phases against a
single experimental apparatus. Phase~1 establishes the contact-force
versus press-depth relationship for each gripper variant at nominal
zero-tilt orientation, and phase~2 perturbs the same measurement with
$\pm 2.5^{\circ}$ tilts about the two orthogonal Cartesian axes.
Together the two phases define a three-dimensional
(depth, tilt-axis, tilt-sign) cell grid whose per-cell measurements
constitute the data corpus reported in Chapter~\ref{ch:results}.

\section{Study Design}
\label{sec:study-design}

The gripper roster comprises six variants organised into two families:
\begin{itemize}
  \item \textbf{Finray family} (four variants):
    \texttt{finray\_18}, \texttt{finray\_18\_model2} (second print of
    the same CAD), \texttt{finray\_26} (larger blade), and
    \texttt{finray\_26\_5deg} (\texttt{finray\_26} on a
    $5^{\circ}$-inclined seat).
  \item \textbf{Parallel-jaw family} (two variants):
    \texttt{parallel\_jaw\_stock} (bare rigid jaws) and
    \texttt{parallel\_jaw\_TPU} (rigid jaws with a $2$\,mm TPU $90$A
    compliant pad).
\end{itemize}

The same fabric specimen (single-layer cotton) rests on the same
acrylic table surface throughout the campaign. The Franka Emika Panda
arm holds the gripper in a fixed base orientation $R_0 = R_x(180^{\circ})$
for phase~1, and applies signed rotations
$R_{\alpha,\beta} = R_x(180^{\circ}) R_x(\alpha) R_y(\beta)$
for phase~2, with $(\alpha, \beta) \in
\{(\pm 2.5^{\circ}, 0), (0, \pm 2.5^{\circ})\}$.

\section{Depth Grid}
\label{sec:depth-grid}

The depth axis is defined relative to the per-session detected cloth
surface via the \texttt{ground\_truth\_finder} touch-off procedure
described in Chapter~\ref{ch:methods}. The depth grid was chosen
adaptively per gripper family:
\begin{itemize}
  \item For the rigid parallel jaws, depths are spaced at $0.25$--$0.5$\,mm
        intervals over $-1$ to $+4$\,mm, close to the narrow
        force-limited working range.
  \item For the finrays, depths are spaced at $0.5$--$1.0$\,mm intervals
        over $-2$ to $+16$\,mm, covering the four regimes
        (boundary, knee, valley, recovery) identified in
        Chapter~\ref{ch:results}.
\end{itemize}

Per (variant, depth, tilt) cell, three trials were collected. Boundary
depths were routinely retested when within-cell standard deviation
exceeded $1$\,N, up to a maximum of two additional cells per
(depth, tilt) combination.

\section{Tilt-Aware Pose Geometry: Verification}
\label{sec:tilt-verify}

Before reporting force results, the implementation of the tilt-aware
tool-centre-point (TCP) tracking equation was verified empirically.
Table~\ref{tab:tilt-geometry} shows the measured flange XY shift during
the stationary grasp sample versus the geometric prediction for
\texttt{finray\_26} at four representative (depth, tilt) cells.

\begin{table}[ht]
\centering
\caption[Tilt geometry verification]{Measured vs.~predicted flange-XY
offset during the stationary grasp sample at four cells. Predicted
offset uses $z_{\text{TCP}} = 94.911$\,mm.}
\label{tab:tilt-geometry}
\small
\begin{tabular}{lccccc}
\toprule
Cell & Tilt cmd & Predicted $\Delta \mathbf{p}_\mathcal{F}$ (mm) &
       Measured $\Delta \mathbf{p}_\mathcal{F}$ (mm) & Residual & Mean $|\mathbf{F}|$ \\
\midrule
$-1.0$\,mm   & $+2.5^{\circ}$ y & $(+4.14, +0.00)$ & $(+4.35, -0.03)$ & $0.21$\,mm & $-8.00$\,N \\
$-1.0$\,mm   & $-2.5^{\circ}$ y & $(-4.14, +0.00)$ & $(-4.02, -0.04)$ & $0.13$\,mm & $-3.39$\,N \\
$+10.0$\,mm  & $+2.5^{\circ}$ y & $(+4.14, +0.00)$ & $(+5.35, -0.28)$ & $1.25$\,mm & $-36.26$\,N \\
$+10.0$\,mm  & $-2.5^{\circ}$ y & $(-4.14, +0.00)$ & $(-2.93, -0.29)$ & $1.21$\,mm & $-35.95$\,N \\
\bottomrule
\end{tabular}
\end{table}

At boundary depth, the residual is $\leq 0.21$\,mm---essentially the
lateral positioning noise of the arm under light load. At safe depth,
both positive and negative tilts exhibit a $+1.2$\,mm common-mode
offset in the same $+x$ direction, regardless of tilt sign. We
interpret this common-mode offset as evidence of finger compliance
under load: the $\sim 36$\,N axial press deflects the soft finray blade
laterally by approximately $1.2$\,mm in a structurally preferred
direction. The differential offset (between $+y$ and $-y$ tilts)
remains correctly predicted by the geometric model in both depth
conditions.

This verification confirms that the control-frame implementation of the
tilt-aware TCP tracker is correct, and any asymmetries observed in the
contact force reflect physical gripper behaviour rather than
control-system errors.

\section{Rotation Tracking Verification}
\label{sec:rotation-verify}

A natural concern is whether the Franka arm faithfully realises the
commanded tilt orientation. The 30\,Hz session log records actual and
commanded RPY at every sample, enabling a direct measurement of
orientation tracking error
\begin{equation}
  \delta(t) = \| \mathrm{RPY}_{\text{actual}}(t) -
                 \mathrm{RPY}_{\text{cmd}}(t) \|_2.
\end{equation}
Across five representative sessions, the mean steady-state $\delta$
was $0.05^{\circ}$ with a worst-case $0.13^{\circ}$. Transient
excursions up to $\sim 5^{\circ}$ occur exclusively within the
\texttt{orientation\_lock} stage at session start, before any descent
or grasp, and are filtered out of trial-level analyses.

A $2.5^{\circ}$ commanded tilt is therefore $50\times$ the noise floor
of the arm's rotation tracking; the observed force differences cannot
be attributed to arm rotation error.

\section{Trial Protocol and Outcome Classification}
\label{sec:trial-protocol}

Each trial follows a fixed six-stage sequence: (1)~set the commanded
pose above the target, (2)~descend to the pre-contact height,
(3)~contact and press to the target depth, (4)~grasp with the gripper's
default close command, (5)~lift by a fixed clearance, (6)~operator
confirms whether the cloth is held. A stationary-sample window of at
least $2$\,seconds during the grasped-and-pressed stage provides the
force measurement.

Trial outcomes are classified into three categories:
\begin{description}
  \item[\textbf{Success}:] the trial executed to completion, contact
        force was sampled cleanly during the stationary window, and the
        cloth remained gripped after the lift stage.
  \item[\textbf{No-grasp}:] the trial executed to completion and
        contact force was sampled cleanly during the stationary window,
        but the cloth slipped from the fingers during or after the lift
        stage. The force measurement remains valid as an
        \emph{engagement-force} characterisation.
  \item[\textbf{Failed trial}:] the trial did not execute to completion.
        The controller either aborted with an excess-force safety event
        or the operator interrupted before the stationary window
        completed. The measurement is invalid and the trial is excluded
        from analyses.
\end{description}

The rationale for this three-way taxonomy---rather than a binary
success/failure classification---is discussed in
Section~\ref{sec:failure-mode-discussion}.
